\documentclass{aastex631}

\usepackage{amsmath,amssymb}
\usepackage{natbib}

\newcommand{\dd}{\mathrm{d}}
\newcommand{\pp}{\partial}
\newcommand{\kbbar}{k_{\rm B}/\mathrm{amu}}
\newcommand{\Mpl}{M_{\rm p}}
\newcommand{\Mcore}{M_{\rm core}}
\newcommand{\Mfe}{M_{\rm Fe}}
\newcommand{\Menv}{M_{\rm env}}
\newcommand{\Rpl}{R_{\rm p}}
\newcommand{\Rcore}{R_{\rm core}}
\newcommand{\Pc}{P_{\rm c}}
\newcommand{\Patm}{P_{\rm atm}}
\newcommand{\Psurf}{P_{\rm surf}}
\newcommand{\Pcb}{P_{\rm cb}}
\newcommand{\orchard}{\textsc{orchard}}

\shorttitle{ORCHARD RK4 Initialization}
\shortauthors{Tejada Arevalo}

\begin{document}

\title{RK4 Shooting Method for Initial Planetary Structure in
  the ORCHARD Evolution Code}

\author{Roberto Tejada Arevalo}

\begin{abstract}
We describe the fourth-order Runge--Kutta (RK4) shooting method used
by the \orchard{} planetary evolution code to construct self-consistent
initial hydrostatic structures at $t=0$.  The method integrates the
equations of hydrostatic equilibrium and mass continuity from the
planetary center to the surface on a fine internal mass grid, then
interpolates the result onto the \orchard{} Lagrangian mesh.  For
envelope-dominated planets (gas giants), a bisection search on the
central pressure $\Pc$ enforces the surface boundary condition
$\Psurf = \Patm$.  For core-dominated planets (super-Earths and
sub-Neptunes with thin envelopes), the surface-pressure shooting is
replaced by an iterative density-matching scheme, because the standard
bisection is ill-conditioned when the envelope mass fraction is small.
We detail the structure equations, the adaptive fine grid construction,
the central boundary condition, the two shooting strategies, and the
interpolation back to the evolution grid.
\end{abstract}

\keywords{planets and satellites: interiors --- planets and satellites:
physical evolution --- methods: numerical}

% ======================================================================
\section{Introduction}\label{sec:intro}
% ======================================================================

The \orchard{} evolution code models the thermal and compositional
evolution of planets on a one-dimensional Lagrangian mass grid of $N$
shells.  At each timestep the transport module updates the entropy and
composition profiles, and a Henyey relaxation solver
(\S\ref{sec:hse_context}) re-establishes hydrostatic equilibrium.  Both
modules require an initial guess for the pressure $P(m)$ and radius
$r(m)$ profiles at $t=0$.

The legacy initialization reads a pre-computed isentrope file (a
tabulated $(r,P)$ structure for a reference gas giant at fixed entropy)
and interpolates it onto the mass grid.  This approach has two
limitations: (i) it is restricted to gas-giant compositions and
entropies for which a tabulated isentrope exists, and (ii) for
non-gas-giant planets (sub-Neptunes, super-Earths) the interpolated
structure can be far from hydrostatic equilibrium, requiring many
initial Henyey iterations to converge.

The RK4 initialization module (\texttt{rk4\_initial.py}) replaces the
isentrope interpolation with a direct numerical integration of the
hydrostatic structure equations.  Given an equation of state, a
composition profile, and a trial central pressure $\Pc$, it integrates
outward in enclosed mass $m$ from the center to the surface using
classical fourth-order Runge--Kutta.  A shooting procedure then adjusts
$\Pc$ until the resulting structure satisfies the desired boundary
condition.

The module supports two shooting strategies
(\S\ref{sec:shooting_gasgiant}, \S\ref{sec:shooting_rocky}), selected
automatically based on the envelope mass fraction.  For gas giants the
standard surface-pressure bisection is used.  For core-dominated planets
a density-matching iteration is employed instead, because the standard
bisection is ill-conditioned when the envelope is a small fraction of
the total mass (\S\ref{sec:why_rocky_fails}).

The remainder of this document is organized as follows.
Section~\ref{sec:equations} presents the structure equations solved by
the RK4 integrator.  Section~\ref{sec:grid} describes the adaptive fine
mass grid.  Section~\ref{sec:bc} details the central boundary condition.
Section~\ref{sec:rk4} gives the RK4 integration scheme.
Section~\ref{sec:shooting_gasgiant} describes the gas-giant shooting
method.  Section~\ref{sec:why_rocky_fails} explains why this fails for
rocky planets.  Section~\ref{sec:shooting_rocky} presents the
density-matching alternative.  Section~\ref{sec:interp} covers the
interpolation onto the \orchard{} grid.  Section~\ref{sec:eos_dispatch}
documents the equation-of-state dispatch across compositional layers.

% ======================================================================
\section{Context: Role of the Initial Structure}\label{sec:hse_context}
% ======================================================================

The RK4 module produces initial profiles $r(m)$, $P(m)$, $\rho(m)$,
and $T(m)$ that serve as the starting point for the evolution.  These
profiles do not need to be exact: the Henyey relaxation solver in
\texttt{hydrostatic.py} will refine them at the first timestep.  What
matters is that the initial guess has:
\begin{enumerate}
  \item a physically reasonable planet radius $\Rpl$,
  \item monotonically decreasing $r(m)$ from surface to center,
  \item monotonically increasing $P(m)$ from surface to center, and
  \item correct equation-of-state dispatch in each compositional layer.
\end{enumerate}
If any of these conditions are violated, the Henyey iteration
diverges and produces NaN values on the first timestep, aborting the
evolution.

% ======================================================================
\section{Structure Equations}\label{sec:equations}
% ======================================================================

The integration solves two coupled ordinary differential equations for
$\ln r$ and $\ln P$ as functions of the enclosed mass $m$, treated as the
independent variable.  Mass continuity gives
\begin{equation}\label{eq:dlnr}
  \frac{\dd \ln r}{\dd m}
  = \frac{1}{4\pi\,r^{3}\,\rho}\,,
\end{equation}
and hydrostatic equilibrium (neglecting rotation) gives
\begin{equation}\label{eq:dlnP}
  \frac{\dd \ln P}{\dd m}
  = -\frac{G\,m}{4\pi\,r^{4}\,P}\,.
\end{equation}
At each point along the integration the density $\rho$ is obtained from
the equation of state as $\rho = \rho(P, S, Y, Z, f_{\rm rock})$, where
$S$ is the entropy, $Y$ the helium mass fraction, $Z$ the heavy-element
mass fraction, and $f_{\rm rock}$ the rock fraction within $Z$.  The
appropriate EOS is selected based on which compositional layer the
current mass coordinate falls in (\S\ref{sec:eos_dispatch}).

The use of logarithmic variables $(\ln r, \ln P)$ rather than $(r, P)$
is essential for numerical stability: both $r$ and $P$ span many orders
of magnitude from center to surface, and the logarithmic form keeps the
dependent variables of order unity.

% ======================================================================
\section{Adaptive Fine Mass Grid}\label{sec:grid}
% ======================================================================

The RK4 integration is performed on a fine internal grid of
$N_{\rm fine} = 2000$ mass points, independent of the $N$-zone
\orchard{} evolution grid.  The fine grid must resolve three distinct
challenges:
\begin{enumerate}
  \item Steep gradients near the planetary center, where $r \propto
    m^{1/3}$ and $\dd r / \dd m$ diverges as $m \to 0$.
  \item The compositional transition at the core--envelope boundary
    $m = \Mcore$, where the EOS switches between mantle/core
    and H/He--metal mixture models.
  \item The thin H/He envelope in core-dominated planets, which can
    contain $< 0.01\%$ of the total mass.
\end{enumerate}

\subsection{Bare Rocky Planets}\label{sec:grid_bare}

For bare rocky planets ($k_{\rm core} = 0$, meaning all $N$ cells are
mantle or iron core with no H/He envelope), the fine grid is simply
$N_{\rm fine} + 1$ logarithmically spaced points from $m_{\rm min}$ to
$\Mpl$.  The log spacing resolves the steep central gradients and
provides uniform fractional resolution throughout the interior.  No
envelope segment is needed.

\subsection{Coreless Planets (Gas Giants)}\label{sec:grid_gasgiant}

For coreless planets ($k_{\rm core} = N$, meaning all cells are
envelope), or when no rocky/icy core is present, the fine grid uses a
two-segment construction:
\begin{itemize}
  \item $N_{\rm log} = 1000$ logarithmically spaced points from
    $m_{\rm min} = 10^{-8}\,\Mpl$ to $m_{\rm split} = 0.5\,\Mpl$, and
  \item $N_{\rm lin} = 1000$ linearly spaced points from $m_{\rm split}$
    to $\Mpl$.
\end{itemize}
The logarithmic spacing resolves the center; the linear spacing gives
uniform resolution in the outer envelope where the atmosphere boundary
condition matters.

\subsection{Core-Dominated Planets}\label{sec:grid_rocky}

When the planet has a rocky/icy core ($0 < k_{\rm core} < N$), the grid
is split at the core--envelope boundary mass $m_{\rm cb} = m_b[k_{\rm
core}]$ rather than at a fixed 50\% of the total mass.  The cell budget
is divided as follows.  Let $f_{\rm env} = (\Mpl - m_{\rm cb}) / \Mpl$
be the envelope mass fraction.  Then
\begin{equation}\label{eq:Nenv}
  N_{\rm env} = \min\!\bigl(\max(50,\;
    \lfloor N_{\rm fine} \cdot f_{\rm env}\rfloor),\;
    N_{\rm fine}/2\bigr),
\end{equation}
and $N_{\rm int} = N_{\rm fine} - N_{\rm env}$.  The minimum of~50
ensures that even an envelope with $f_{\rm env} \sim 10^{-4}$ receives
at least 50~fine cells, providing adequate support for the subsequent
interpolation back onto the $N$-zone evolution grid
(\S\ref{sec:interp}).  The two segments are:
\begin{itemize}
  \item \emph{Interior} ($N_{\rm int}$ points): logarithmically spaced
    from $m_{\rm min}$ to $m_{\rm cb}$.  This covers the iron core and
    silicate mantle, with the logarithmic spacing concentrating
    resolution near the center.
  \item \emph{Envelope} ($N_{\rm env}$ points): logarithmically spaced
    from $m_{\rm cb}$ to $\Mpl$.  The log spacing concentrates points
    near the core--envelope boundary (where the EOS transition occurs)
    and near the surface.
\end{itemize}
The core boundary mass $m_{\rm cb}$ is an exact grid point, ensuring
that the RK4 integration steps cleanly through the EOS transition rather
than straddling it.

\paragraph{Why the fixed grid fails for rocky planets.}
With the original two-segment grid (\S\ref{sec:grid_gasgiant}), the
linear outer segment has a step size of $\Delta m \approx 0.5\,\Mpl /
N_{\rm lin} = 5 \times 10^{-4}\,\Mpl$.  For a super-Earth with
$f_{\rm env} = 10^{-4}$, the entire envelope falls within a single step
of this grid.  No fine-grid points sample the envelope EOS, and the
subsequent interpolation (\S\ref{sec:interp}) extrapolates flat $r$ and
$P$ values into all envelope cells of the \orchard{} grid, producing
non-monotonic profiles that cause the Henyey solver to diverge.

% ======================================================================
\section{Central Boundary Condition}\label{sec:bc}
% ======================================================================

The integration starts at the innermost fine-grid point $m_0 = m_{\rm
min} = 10^{-8}\,\Mpl$ with boundary conditions derived from the
uniform-density-sphere regularity condition.  Given a trial central
pressure $\Pc$ and the central density $\rho_{\rm c} =
\rho(\Pc,\,S_{\rm c},\,Y_{\rm c},\,Z_{\rm c})$ from the EOS at the
center:
\begin{equation}\label{eq:r_center}
  r_0 = \left(\frac{3\,m_0}{4\pi\,\rho_{\rm c}}\right)^{1/3},
\end{equation}
\begin{equation}\label{eq:P_center}
  P_0 = \Pc - \frac{3G}{8\pi}
    \left(\frac{4\pi\,\rho_{\rm c}}{3}\right)^{4/3}\!m_0^{2/3}\,.
\end{equation}
The correction term in Equation~\eqref{eq:P_center} accounts for the
pressure drop across the innermost sphere of mass $m_0$, ensuring that
the integration starts from a self-consistent state rather than imposing
$P = \Pc$ at $m = m_0 > 0$.  If the correction would reduce $P_0$ below
$\Pc/2$ (which can happen for extremely low $\Pc$ guesses), it is
clamped to $P_0 = \Pc/2$ to prevent numerical underflow.

% ======================================================================
\section{RK4 Integration}\label{sec:rk4}
% ======================================================================

The integration advances from the center ($j=0$) to the surface
($j = N_{\rm fine}$) using the classical fourth-order Runge--Kutta
scheme.  At each step, the mass increment is $h_j = m_{j+1} - m_j$ and
the state vector is $\mathbf{u}_j = (\ln r_j,\;\ln P_j)$.  The four
RK4 stages are:
\begin{equation}\label{eq:rk4_stages}
\begin{aligned}
  \mathbf{k}_1 &= \mathbf{f}(m_j,\;\mathbf{u}_j), \\
  \mathbf{k}_2 &= \mathbf{f}\!\left(m_j + \tfrac{h}{2},\;
    \mathbf{u}_j + \tfrac{h}{2}\,\mathbf{k}_1\right), \\
  \mathbf{k}_3 &= \mathbf{f}\!\left(m_j + \tfrac{h}{2},\;
    \mathbf{u}_j + \tfrac{h}{2}\,\mathbf{k}_2\right), \\
  \mathbf{k}_4 &= \mathbf{f}\!\left(m_{j+1},\;
    \mathbf{u}_j + h\,\mathbf{k}_3\right), \\
  \mathbf{u}_{j+1} &= \mathbf{u}_j + \frac{h}{6}\bigl(
    \mathbf{k}_1 + 2\,\mathbf{k}_2 + 2\,\mathbf{k}_3
    + \mathbf{k}_4\bigr),
\end{aligned}
\end{equation}
where $\mathbf{f} = (\dd\ln r/\dd m,\;\dd\ln P/\dd m)$ from
Equations~\eqref{eq:dlnr}--\eqref{eq:dlnP}.  The density $\rho$ is
re-evaluated from the EOS at each sub-step using the current (trial)
pressure.  For the stages at the midpoint $m_j + h/2$, the composition
$(S, Y, Z, f_{\rm rock})$ is taken from the fine-grid cell $j$; for the
final stage at $m_{j+1}$, the composition is taken from cell
$\min(j+1,\,N_{\rm fine}-1)$.

Two safety measures prevent the integration from producing unphysical
results:
\begin{enumerate}
  \item \textbf{Density floor.}  If the EOS returns $\rho \leq 0$ or a
    non-finite value, a floor value $\rho_{\rm floor} = 10^{-6}\;\mathrm
    {g\;cm^{-3}}$ is substituted.  This prevents $\dd\ln r/\dd m$ from
    diverging in the low-density outer envelope.
  \item \textbf{Radius cap.}  The radius is capped at $r_{\rm max} =
    10^{12}\;\mathrm{cm}$ ($\approx 100\;R_{\rm J}$) to prevent
    floating-point overflow in $\exp(3\ln r)$.  If either $\ln r$ or
    $\ln P$ becomes non-finite after an RK4 step, the state is frozen at
    the last valid values.
\end{enumerate}

% ======================================================================
\section{Shooting: Gas-Giant Mode}\label{sec:shooting_gasgiant}
% ======================================================================

For envelope-dominated planets where the core mass fraction is less than
90\% of the total, the central pressure is determined by a standard
bisection shooting method.  The target is to match the surface pressure:
\begin{equation}\label{eq:shooting_target}
  P_{\rm surf}(\Pc) = \Patm\,,
\end{equation}
where $\Patm$ is the atmospheric pressure from the parameter file
(typically $10^6\;\mathrm{dyn\;cm^{-2}} = 1\;\mathrm{bar}$).

\subsection{Initial Guess and Bracketing}

The initial guess for $\Pc$ is based on the planet-to-Jupiter mass ratio:
\begin{equation}\label{eq:Pc_guess}
  \Pc^{(0)} = 4 \times 10^{13}
    \cdot\bigl[\max(M_{\rm p}/M_{\rm J},\;0.1)\bigr]^{2}
    \;\mathrm{dyn\;cm^{-2}}\,,
\end{equation}
clipped to the range $[10^{11},\,10^{16}]\;\mathrm{dyn\;cm^{-2}}$.
The bisection brackets are initialized at $P_{\rm lo} = 0.01\,\Pc^{(0)}$
and $P_{\rm hi} = 100\,\Pc^{(0)}$.  If the initial brackets do not
straddle the root, they are expanded by factors of~10 up to 10~times.

\subsection{Bisection Iteration}

At each iteration, the geometric midpoint $\Pc = \sqrt{P_{\rm lo}\,
P_{\rm hi}}$ is evaluated by running the full RK4 integration and
reading off $\Psurf = P_{\rm cells}[0]$.  If $\Psurf > \Patm$ the
central pressure is too high, and $P_{\rm hi} \gets \Pc$; otherwise
$P_{\rm lo} \gets \Pc$.  Convergence is declared when the relative error
$|\Psurf - \Patm|/\Patm < 1\%$.  The maximum number of iterations
is~60.

% ======================================================================
\section{Why Gas-Giant Shooting Fails for Rocky
  Planets}\label{sec:why_rocky_fails}
% ======================================================================

For core-dominated planets with thin H/He envelopes, the surface-pressure
bisection (\S\ref{sec:shooting_gasgiant}) fails for two reasons:

\paragraph{(1) The surface pressure function is non-monotonic.}
The function $\Psurf(\Pc)$ is well-behaved and monotonically increasing
for gas giants: increasing $\Pc$ compresses the interior more and raises
the surface pressure.  For rocky planets with thin envelopes, however,
$\Psurf(\Pc)$ is wildly non-monotonic.  At some values of $\Pc$ the
pressure underflows to zero partway through the integration; at others
the integration completes but with $\Psurf$ orders of magnitude above
$\Patm$.  This non-monotonicity arises because the thin envelope carries
negligible mass and therefore exerts negligible gravitational influence
on the pressure profile.  The surface pressure is determined almost
entirely by the pressure at the core--envelope boundary, which itself
depends sensitively on the interior structure.  The bisection can
converge to a spurious root where $\Psurf \approx \Patm$ by accident,
but the resulting structure has an unphysically small planet radius and
non-monotonic profiles.

\paragraph{(2) No physical root exists.}
For a physically correct central pressure (e.g., $\Pc \approx 3.6
\times 10^{12}\;\mathrm{dyn\;cm^{-2}}$ for a $1\;M_\oplus$ planet), the
mantle integration produces a core-boundary pressure $\Pcb \sim
10^{9}$--$10^{12}\;\mathrm{dyn\;cm^{-2}}$.  The thin envelope (mass
fraction $\sim 10^{-4}$) cannot reduce $P$ from $\Pcb$ to
$\Patm = 10^{6}\;\mathrm{dyn\;cm^{-2}}$.  The fractional pressure change
across the envelope is
\begin{equation}\label{eq:DeltalnP}
  \Delta\ln P \approx
    -\frac{G\,\Mpl}{4\pi\,\Rcore^{4}\,\Pcb}\;\Menv\,,
\end{equation}
which for $\Mpl = M_\oplus$, $\Rcore \approx R_\oplus$, $\Pcb = 10^{9}$,
and $\Menv = 10^{-4}\,\Mpl$ evaluates to $|\Delta\ln P| \sim 0.1$,
i.e., a $\sim 10\%$ pressure drop---far short of the factor of
$\sim 10^{3}$--$10^{6}$ needed.

In contrast, the Henyey relaxation solver in the main evolution code
handles this correctly because it imposes $P[0] = \Patm$ as a fixed
outer boundary condition and solves inward, rather than integrating
outward from the center.

% ======================================================================
\section{Shooting: Rocky-Planet Mode}\label{sec:shooting_rocky}
% ======================================================================

For planets where the envelope mass fraction is less than~10\% of the
total mass, the module switches to a density-matching iteration that
seeks a physically reasonable interior structure rather than matching
the surface pressure.  The surface pressure will be corrected by the
Henyey relaxation solver at the first evolution timestep.

\subsection{Classification}\label{sec:classification}

A planet is classified as ``rocky'' (core-dominated) in two cases:
\begin{enumerate}
  \item \textbf{Bare rock} ($k_{\rm core} = 0$): all cells are
    mantle/core with no H/He envelope.  This corresponds to
    $M_{\rm core} = \Mpl$ in the parameter file.
  \item \textbf{Thin envelope} ($0 < k_{\rm core} < N$ and
    $f_{\rm env} < 0.1$): the envelope mass fraction
    \begin{equation}\label{eq:rocky_criterion}
      f_{\rm env} \equiv \frac{\Mpl - m_b[k_{\rm core}]}{\Mpl}
    \end{equation}
    is less than 10\%.
\end{enumerate}
Together these cases capture bare super-Earths, super-Earths with thin
atmospheres, and sub-Neptunes with small volatile envelopes.  Gas giants
and ice giants with substantial envelopes ($k_{\rm core} \geq N$ or
$f_{\rm env} \geq 0.1$) use the standard bisection
(\S\ref{sec:shooting_gasgiant}).

\subsection{Central Pressure Estimate}

The initial estimate of $\Pc$ uses the uniform-sphere formula.  Let
$\bar{\rho}$ be a characteristic mean density, obtained from the mantle
EOS evaluated at the mantle entropy $S_{\rm mantle}$ and a reference
pressure of $100\;\mathrm{GPa}$:
\begin{equation}\label{eq:rhobar}
  \bar{\rho} = \rho_{\rm mantle}(S_{\rm mantle},\;100\;\mathrm{GPa})\,.
\end{equation}
The uniform-sphere central pressure is then
\begin{equation}\label{eq:Pc_est}
  \Pc^{(0)} = \frac{3G}{8\pi}
    \left(\frac{4\pi\,\bar{\rho}}{3}\right)^{4/3}\!\Mpl^{2/3}\,,
\end{equation}
clipped to $[10^{10},\,10^{16}]\;\mathrm{dyn\;cm^{-2}}$.  The target
planet radius is similarly estimated as
\begin{equation}\label{eq:R_target}
  R_{\rm target} = \left(\frac{3\,\Mpl}
    {4\pi\,\bar{\rho}}\right)^{1/3}\!.
\end{equation}

\subsection{Iterative Radius Matching}

Starting from $\Pc^{(0)}$, the iteration proceeds as follows:
\begin{enumerate}
  \item Perform the full RK4 integration at the current $\Pc$ to obtain
    $r_b$, $P_{\rm cells}$, $\rho_{\rm cells}$, and $T_{\rm cells}$.
  \item Evaluate a quality score:
    \begin{equation}\label{eq:score}
      \mathcal{S} = \bigl|\ln(\Rpl / R_{\rm target})\bigr|
        + \begin{cases} 10 & \text{if profiles are non-monotonic,} \\
                         0  & \text{otherwise.}\end{cases}
    \end{equation}
  \item If the profiles are monotonic and $|\Rpl / R_{\rm target} - 1|
    < 0.3$, accept the solution.
  \item Otherwise, adjust $\Pc$ using a damped correction:
    \begin{equation}\label{eq:Pc_update}
      \Pc \gets \Pc \times
        \mathrm{clip}\!\left[\left(\frac{\Rpl}{R_{\rm target}}
        \right)^{-2},\;0.3,\;3.0\right].
    \end{equation}
    The exponent $-2$ is a damped version of the $\Rpl^{-4}$ scaling
    expected from hydrostatic equilibrium ($\Pc \propto \Mpl^{2} /
    \Rpl^{4}$), and the clipping prevents overshooting.
  \item Repeat up to 30 iterations, tracking the best result (lowest
    $\mathcal{S}$).
\end{enumerate}

In practice, convergence typically occurs within 2--5 iterations.  For
the benchmark case of a $1\;M_\oplus$ super-Earth with $M_{\rm core} =
0.9999\;M_\oplus$, the iteration converges in 3~iterations to $\Pc =
4.3 \times 10^{12}\;\mathrm{dyn\;cm^{-2}}$ and $\Rpl = 1.03\;R_\oplus$.

% ======================================================================
\section{Interpolation onto the Evolution Grid}\label{sec:interp}
% ======================================================================

After the RK4 integration, the fine-grid profiles $\ln r(m)$ and
$\ln P(m)$ are interpolated onto the \orchard{} mass boundaries
$m_b[k]$ ($k = 0,\ldots,N$, descending from $\Mpl$ at the surface to
$\sim 0$ at the center).

\subsection{Boundary Values}

Linear interpolation is performed in $(\ln r,\,m)$ and $(\ln P,\,m)$
space using \texttt{scipy.interpolate.interp1d}.  The mass coordinates
are clamped to $[m_{\rm min},\,\Mpl]$ to avoid extrapolation at the
center where $m_b[N] \approx 0 < m_{\rm min}$.  The boundary radius and
pressure are recovered as
\begin{equation}
  r_b[k] = \exp\!\bigl(\hat{f}_{\ln r}(m_b[k])\bigr), \quad
  P_b[k] = \exp\!\bigl(\hat{f}_{\ln P}(m_b[k])\bigr),
\end{equation}
where $\hat{f}$ denotes the interpolant.

\subsection{Cell-Centered Values}

Cell-centered pressures are computed as the geometric mean of the
bounding pressures:
\begin{equation}\label{eq:Pcell}
  P_k = \sqrt{P_b[k]\,P_b[k+1]}\,,
  \qquad k = 0,\ldots,N-1\,.
\end{equation}
The innermost cell ($k = N-1$) is clamped to $P_{N-1} \geq \Pc / 2$ to
prevent underflow from the extrapolation at $m \approx 0$.

\subsection{Density and Temperature Recovery}

With $P_k$ determined, the density and temperature at each cell center
are recomputed from the equation of state appropriate to that cell's
compositional layer (\S\ref{sec:eos_dispatch}).  This ensures
thermodynamic consistency between the pressure, density, and temperature
profiles even if the EOS evaluations during the RK4 integration were
approximate (due to the finite step size).

% ======================================================================
\section{Equation-of-State Dispatch}\label{sec:eos_dispatch}
% ======================================================================

Both the fine-grid RK4 integration and the final cell-centered density
recovery use the same EOS dispatch based on the cell's compositional
layer.  The \orchard{} grid is divided into three regions by the indices
$k_{\rm core}$ and $k_{\rm Fe}$:

\begin{enumerate}
  \item \textbf{Envelope} ($k < k_{\rm core}$):
    Hydrogen--helium--metal mixture EOS.  Given entropy $S$, $\log_{10}
    P$, helium fraction $Y$, metallicity $Z$, and rock fraction
    $f_{\rm rock}$, returns $\log_{10}\rho$ and $\log_{10}T$.  Available
    backends: Chabrier \& Debras (2021, ``cd''), Chabrier et al.\
    (2019, ``cms''), or Saumon et al.\ (1995, ``scvh'').

  \item \textbf{Mantle} ($k_{\rm core} \leq k < k_{\rm Fe}$):
    Silicate EOS (Mg$_2$SiO$_4$ ANEOS, MgSiO$_3$, or AQUA water).
    Given entropy $S_{\rm mantle}$ and pressure $P$ in GPa, returns
    $\rho$ and $T$.

  \item \textbf{Iron core} ($k \geq k_{\rm Fe}$):
    Iron alloy or pure iron EOS.  Given entropy $S_{\rm core}$ and $P$
    in GPa, returns $T$ via $(s, P)$ inversion, then $\rho$ via $(P,T)$
    lookup.
\end{enumerate}

During the RK4 integration, the composition at each fine-grid cell is
looked up from the \orchard{} grid via a \texttt{searchsorted} mapping
from the fine-grid mass coordinates to the \orchard{} cell indices.
This ensures that the EOS transition occurs at the correct mass
coordinate regardless of the fine-grid spacing.

For the envelope, composition is set to $S = S_{\rm ini}$, $Y = Y_{\rm
ini}\,(1 - Z_{\rm ini})$, $Z = Z_{\rm ini}$, and $f_{\rm rock} =
f_{\rm rock,\,ini}$.  For mantle and core cells, $Y = Z = f_{\rm rock}
= 0$ and the entropy is set to $S_{\rm mantle}$ and $S_{\rm core}$
respectively.

% ======================================================================
\section{Summary}\label{sec:summary}
% ======================================================================

The RK4 initialization module provides \orchard{} with self-consistent
initial hydrostatic structures for any planet mass, composition, and
initial entropy, replacing the legacy approach of interpolating a
pre-computed gas-giant isentrope.  Two key design choices enable it to
work across the full range of planet types from gas giants to bare
super-Earths:

\begin{enumerate}
  \item \textbf{Adaptive fine grid} (\S\ref{sec:grid}): The fine mass
    grid adapts to the planet architecture.  Bare rocky planets use a
    single log-spaced segment (\S\ref{sec:grid_bare}).  Planets with
    cores and envelopes split the grid at the core--envelope boundary,
    guaranteeing a minimum of 50~fine cells in the envelope even when
    it contains $< 0.01\%$ of the total mass
    (\S\ref{sec:grid_rocky}).  Gas giants use a two-segment log/linear
    construction (\S\ref{sec:grid_gasgiant}).  This prevents the
    interpolation from producing flat, non-monotonic profiles.

  \item \textbf{Dual shooting strategy} (\S\ref{sec:shooting_gasgiant},
    \S\ref{sec:shooting_rocky}): Gas giants use standard bisection on
    $\Psurf = \Patm$.  Core-dominated planets---both bare rocks
    ($k_{\rm core} = 0$) and thin-envelope sub-Neptunes/super-Earths---use
    iterative density matching to find a $\Pc$ that produces monotonic
    profiles with a physically reasonable radius, since the surface-%
    pressure bisection is ill-conditioned when the envelope mass fraction
    is small or zero (\S\ref{sec:why_rocky_fails}).  In both cases, the
    Henyey relaxation solver refines the structure at the first evolution
    timestep.
\end{enumerate}

\end{document}
